\chapter{Title of chapter 2}

The title of a subsequent chapter should always appear on a new page. Every paragraph containing the title of a chapter or subchapter has the following intervals:
\begin{itemize}
	\item top 12 pt., 
	\item bottom 6 pt.
\end{itemize}

Here is some text to get rid of Underfull Badbox. Here is some text to get rid of Underfull Badbox. Here is some text to get rid of Underfull Badbox. Here is some text to get rid of Underfull Badbox. Here is some text to get rid of Underfull Badbox. Here is some text to get rid of Underfull Badbox. Here is some text to get rid of Underfull Badbox. Here is some text to get rid of Underfull Badbox. Here is some text to get rid of Underfull Badbox. Here is some text to get rid of Underfull Badbox. Here is some text to get rid of Underfull Badbox. Here is some text to get rid of Underfull Badbox. Here is some text to get rid of Underfull Badbox. Here is some text to get rid of Underfull Badbox. Here is some text to get rid of Underfull Badbox. Here is some text to get rid of Underfull Badbox. Here is some text to get rid of Underfull Badbox.



\section{Subchapter title}

References in the text to literature should show in square brackets their reference number from the \textbf{Bibliography} \cite{html5-perf} or the name of the author/s and year of publication [Kowalski J., 2002] and page number in the case of citations [Kowalski J., 2002, p. 3].

The \textbf{Bibliography} should be arranged alphabetically or in the order items are mentioned in the chapter. If several items of literature are referred to at one point in the text, a dash may be used in the square brackets [1-5] if they are in consecutive order in the \textbf{Bibliography}, or if their order is not consecutive, using commas, as follows: [1, 3, 5]. 

An example of bibliographic entries is shown in the Bibliography chapter.

Figures are numbered in the order in which they appear in the chapter. In the caption below a figure, the number of the figure is preceded by the abbreviation ‘Fig.’ and the appropriate chapter number, e.g. ‘Fig.~2.1.’

Figures appearing in the thesis should be centred. The spacing in a paragraph containing a figure should be:
\begin{itemize}
	\item top 12 pt,
	\item bottom 0 pt.
\end{itemize}

Every figure must be referred to in the text, for example as follows: `Fig.~\ref{fig:ETI-logo}. presents\ldots'. 

If the thesis is written in English, all the tables and figures must also be presented in the English language.

An example of a properly presented figure and caption is provided below.

\begin{figure}[h]
	\includegraphics{./img/LogoETI_en}
	\caption{Emblem of Gdańsk University of Technology.}
	\label{fig:ETI-logo}
\end{figure}

The caption below the figure should be centered, have a 9 pt. font size and end with a full stop. The top margin should be 6 pt., the bottom margin 12 pt. and the spacing should be single.

The first paragraph below a figure should have a top margin of 12 pt.



\subsection{Subchapter section heading}

Every paragraph containing a subchapter section heading should have the following spacing:
\begin{itemize}
	\item top 12 pt., 
	\item bottom 6 pt.
\end{itemize}



\section{Subchapter heading}

Variables should be written in italics, e.g. $x, ni, ni+1$, whereas symbols denoting vectors or matrices should be written in bold, e.g. $\mathbf{v}, \mathbf{A}$. The minus sign should directly precede the number without a space, e.g. $-20$. There must always be a space between a numerical value and unit of measurement, e.g. $1\,\mathrm{V}$ or $10\,\mathrm{km}$.

Preferably, the insertion of equations (INSERT -\textgreater{} Equation) directly into the text should be avoided. If there is a possibility of legibly inserting an equation in one line, it should be done as in the following examples: $1/2t^2$ or $e^{2x+1}$. This should be immediately followed by an explanation as to what the particular symbols denote, e.g. where: $t$ --– time [s].

Equations that might become illegible when written in a single line should be written in a separate paragraph, e.g.
\begin{equation}
	\label{eq:exemplary-equation}
	s = v_0 \cdot t + \frac{a \cdot t^2}{2}
\end{equation}
where: \\
\begin{tabularx}{\textwidth}{p{1.25cm}p{0.5cm}X}
	$s$		& ---	& linear displacement with constant acceleration [m], \\
	$v_0$ 	& ---	& initial velocity [m/s], \\
	$t$		& ---	& time of object motion [s], \\
	$a$		& ---	& acceleration [m/s\textsuperscript{2}].
\end{tabularx}

Every equation should be centered. Its number should be preceded by the chapter number and a dot, in parentheses. Every equation has to be referred to in the text, e.g. `Equation~\eqref{eq:exemplary-equation} allows us to estimate\ldots'.

Any footnotes should appear below a line at the bottom of the page,\footnote{Footnotes should be written using Arial 9 pt.} and there numbering should be consecutive throughout the thesis. To make footnotes use NUMBERING OPTIONS -\textgreater{} add footnote.